% Options for packages loaded elsewhere
% Options for packages loaded elsewhere
\PassOptionsToPackage{unicode}{hyperref}
\PassOptionsToPackage{hyphens}{url}
\PassOptionsToPackage{dvipsnames,svgnames,x11names}{xcolor}
%
\documentclass[
  letterpaper,
  DIV=11,
  numbers=noendperiod]{scrartcl}
\usepackage{xcolor}
\usepackage{amsmath,amssymb}
\setcounter{secnumdepth}{5}
\usepackage{iftex}
\ifPDFTeX
  \usepackage[T1]{fontenc}
  \usepackage[utf8]{inputenc}
  \usepackage{textcomp} % provide euro and other symbols
\else % if luatex or xetex
  \usepackage{unicode-math} % this also loads fontspec
  \defaultfontfeatures{Scale=MatchLowercase}
  \defaultfontfeatures[\rmfamily]{Ligatures=TeX,Scale=1}
\fi
\usepackage{lmodern}
\ifPDFTeX\else
  % xetex/luatex font selection
  \setmainfont[]{Noto Color Emoji}
\fi
% Use upquote if available, for straight quotes in verbatim environments
\IfFileExists{upquote.sty}{\usepackage{upquote}}{}
\IfFileExists{microtype.sty}{% use microtype if available
  \usepackage[]{microtype}
  \UseMicrotypeSet[protrusion]{basicmath} % disable protrusion for tt fonts
}{}
\makeatletter
\@ifundefined{KOMAClassName}{% if non-KOMA class
  \IfFileExists{parskip.sty}{%
    \usepackage{parskip}
  }{% else
    \setlength{\parindent}{0pt}
    \setlength{\parskip}{6pt plus 2pt minus 1pt}}
}{% if KOMA class
  \KOMAoptions{parskip=half}}
\makeatother
% Make \paragraph and \subparagraph free-standing
\makeatletter
\ifx\paragraph\undefined\else
  \let\oldparagraph\paragraph
  \renewcommand{\paragraph}{
    \@ifstar
      \xxxParagraphStar
      \xxxParagraphNoStar
  }
  \newcommand{\xxxParagraphStar}[1]{\oldparagraph*{#1}\mbox{}}
  \newcommand{\xxxParagraphNoStar}[1]{\oldparagraph{#1}\mbox{}}
\fi
\ifx\subparagraph\undefined\else
  \let\oldsubparagraph\subparagraph
  \renewcommand{\subparagraph}{
    \@ifstar
      \xxxSubParagraphStar
      \xxxSubParagraphNoStar
  }
  \newcommand{\xxxSubParagraphStar}[1]{\oldsubparagraph*{#1}\mbox{}}
  \newcommand{\xxxSubParagraphNoStar}[1]{\oldsubparagraph{#1}\mbox{}}
\fi
\makeatother


\usepackage{longtable,booktabs,array}
\newcounter{none} % for unnumbered tables
\usepackage{calc} % for calculating minipage widths
% Correct order of tables after \paragraph or \subparagraph
\usepackage{etoolbox}
\makeatletter
\patchcmd\longtable{\par}{\if@noskipsec\mbox{}\fi\par}{}{}
\makeatother
% Allow footnotes in longtable head/foot
\IfFileExists{footnotehyper.sty}{\usepackage{footnotehyper}}{\usepackage{footnote}}
\makesavenoteenv{longtable}
\usepackage{graphicx}
\makeatletter
\newsavebox\pandoc@box
\newcommand*\pandocbounded[1]{% scales image to fit in text height/width
  \sbox\pandoc@box{#1}%
  \Gscale@div\@tempa{\textheight}{\dimexpr\ht\pandoc@box+\dp\pandoc@box\relax}%
  \Gscale@div\@tempb{\linewidth}{\wd\pandoc@box}%
  \ifdim\@tempb\p@<\@tempa\p@\let\@tempa\@tempb\fi% select the smaller of both
  \ifdim\@tempa\p@<\p@\scalebox{\@tempa}{\usebox\pandoc@box}%
  \else\usebox{\pandoc@box}%
  \fi%
}
% Set default figure placement to htbp
\def\fps@figure{htbp}
\makeatother


% definitions for citeproc citations
\NewDocumentCommand\citeproctext{}{}
\NewDocumentCommand\citeproc{mm}{%
  \begingroup\def\citeproctext{#2}\cite{#1}\endgroup}
\makeatletter
 % allow citations to break across lines
 \let\@cite@ofmt\@firstofone
 % avoid brackets around text for \cite:
 \def\@biblabel#1{}
 \def\@cite#1#2{{#1\if@tempswa , #2\fi}}
\makeatother
\newlength{\cslhangindent}
\setlength{\cslhangindent}{1.5em}
\newlength{\csllabelwidth}
\setlength{\csllabelwidth}{3em}
\newenvironment{CSLReferences}[2] % #1 hanging-indent, #2 entry-spacing
 {\begin{list}{}{%
  \setlength{\itemindent}{0pt}
  \setlength{\leftmargin}{0pt}
  \setlength{\parsep}{0pt}
  % turn on hanging indent if param 1 is 1
  \ifodd #1
   \setlength{\leftmargin}{\cslhangindent}
   \setlength{\itemindent}{-1\cslhangindent}
  \fi
  % set entry spacing
  \setlength{\itemsep}{#2\baselineskip}}}
 {\end{list}}
\usepackage{calc}
\newcommand{\CSLBlock}[1]{\hfill\break\parbox[t]{\linewidth}{\strut\ignorespaces#1\strut}}
\newcommand{\CSLLeftMargin}[1]{\parbox[t]{\csllabelwidth}{\strut#1\strut}}
\newcommand{\CSLRightInline}[1]{\parbox[t]{\linewidth - \csllabelwidth}{\strut#1\strut}}
\newcommand{\CSLIndent}[1]{\hspace{\cslhangindent}#1}



\setlength{\emergencystretch}{3em} % prevent overfull lines

\providecommand{\tightlist}{%
  \setlength{\itemsep}{0pt}\setlength{\parskip}{0pt}}



 


\KOMAoption{captions}{tableheading}
\makeatletter
\@ifpackageloaded{caption}{}{\usepackage{caption}}
\AtBeginDocument{%
\ifdefined\contentsname
  \renewcommand*\contentsname{Table of contents}
\else
  \newcommand\contentsname{Table of contents}
\fi
\ifdefined\listfigurename
  \renewcommand*\listfigurename{List of Figures}
\else
  \newcommand\listfigurename{List of Figures}
\fi
\ifdefined\listtablename
  \renewcommand*\listtablename{List of Tables}
\else
  \newcommand\listtablename{List of Tables}
\fi
\ifdefined\figurename
  \renewcommand*\figurename{Figure}
\else
  \newcommand\figurename{Figure}
\fi
\ifdefined\tablename
  \renewcommand*\tablename{Table}
\else
  \newcommand\tablename{Table}
\fi
}
\@ifpackageloaded{float}{}{\usepackage{float}}
\floatstyle{ruled}
\@ifundefined{c@chapter}{\newfloat{codelisting}{h}{lop}}{\newfloat{codelisting}{h}{lop}[chapter]}
\floatname{codelisting}{Listing}
\newcommand*\listoflistings{\listof{codelisting}{List of Listings}}
\makeatother
\makeatletter
\makeatother
\makeatletter
\@ifpackageloaded{caption}{}{\usepackage{caption}}
\@ifpackageloaded{subcaption}{}{\usepackage{subcaption}}
\makeatother
\usepackage{bookmark}
\IfFileExists{xurl.sty}{\usepackage{xurl}}{} % add URL line breaks if available
\urlstyle{same}
% fallback for those not using the hyperref driver hyperxmp:
\makeatletter
\@ifundefined{xmpquote}{\newcommand{\xmpquote}[1]{#1}}{}
\makeatother
\hypersetup{
  pdftitle={CÁLCULO III - A RCN00067 -- 2026-2},
  pdfauthor={Prof.~Reginaldo Demarque},
  colorlinks=true,
  linkcolor={blue},
  filecolor={Maroon},
  citecolor={Blue},
  urlcolor={Blue},
  pdfcreator={LaTeX via pandoc}}


\title{\href{https://demarque.mat.br/calculoIII/2026-1/c3-A-2026-1.html}{CÁLCULO
III - A RCN00067 -- 2026-2}}
\author{\href{https://www.demarque.mat.br}{Prof.~Reginaldo Demarque}}
\date{}
\begin{document}
\maketitle

\renewcommand*\contentsname{Table of contents}
{
\setcounter{tocdepth}{3}
\tableofcontents
}

\section{Informações Básicas}\label{informauxe7uxf5es-buxe1sicas}

\begin{itemize}
\item
  \textbf{Início e Fim do Período:} 03/08/2026 a 11/12/2026.
\item
  \href{https://demarque.mat.br/calculoIII/2026-2/c3-A-2026-2.html}{Página
  da disciplina}
\end{itemize}

{\def\LTcaptype{none} % do not increment counter
\begin{longtable}[]{@{}lll@{}}
\toprule\noalign{}
Dia & Horário & Sala \\
\midrule\noalign{}
\endhead
\bottomrule\noalign{}
\endlastfoot
Segunda & 9h - 11h & Sala S01 IHS Prédio Alugado \\
Terça & 9h - 11h & Sala S01 IHS Prédio Alugado \\
\end{longtable}
}

\begin{itemize}
\item
  \textbf{Material a ser utilizado}

  \begin{itemize}
  \item
    \href{https://reginaldodr.github.io/slides/calculo3/calculo3.pdf}{slides}
  \item
    \href{https://demarque.mat.br/calculoIII/provas/}{VAs Antigas}
  \end{itemize}
\item
  \textbf{Ementa:} Busque por RCN00075 em
  \href{https://app.uff.br/graduacao/quadrodehorarios/}{quadro de
  horários}:\\
  EQUAÇÕES PARAMÉTRICAS. FUNÇÕES VETORIAIS. SUPERFÍCIES QUÁDRICAS E
  CILÍNDRICAS. FUNÇÕES DE VÁRIAS VARIÁVEIS, LIMITE, CONTINUIDADE,
  DERIVADAS DIRECIONAIS, OTIMIZAÇÃO E MULTIPLICADORES DE LAGRANGE.
\end{itemize}

\section{Apresentação do Curso}\label{apresentauxe7uxe3o-do-curso}

Este curso é constituído de dois módulos: \textbf{Funções Vetoriais e
Superfícies} e \textbf{Funções de Várias Variáveis}.

O \textbf{livro texto} será referência {[}1{]}, que \textbf{está
disponíveis para empréstimo na Biblioteca do campus} --
\href{https://bibliotecas.uff.br/bro/}{BRO}. É altamente recomendado que
vocês consultem este material, de preferência o livro físico.

As demais referência contidas na Section~\ref{sec-ref} formam uma lista
complementar/auxiliar do curso, foram usados principalmente na
elaboração das notas de aulas e exemplos.

\section{Avaliação}\label{avaliauxe7uxe3o}

A avaliação será feita mediante \textbf{2 Verificações de Aprendizagem
(VAs) Escritas}. A média parial (NP) será a média aritmética entre elas,
isto é, \[NP=\frac{VA_1+VA_2}{2}\] De acordo com o Regulamento da UFF, a
\textbf{Nota Final} é calculada da seguinte forma:

\begin{quote}
Definida como sendo igual à média parcial, caso o discente tenha obtido
aprovação direta, ou igual a 6(seis), se a aprovação foi obtida na
verificação suplementar (VS). No caso de reprovação na VS, a nota final
será o resultado do cálculo da média aritmética entre a média parcial e
a nota obtida na VS.
\end{quote}

\subsection*{Objetivos}\label{objetivos}
\addcontentsline{toc}{subsection}{Objetivos}

As VAs tem como objetivos avaliar se o aluno:

\begin{itemize}
\tightlist
\item
  É capaz de interpretar as questões corretamente e formular soluções
  para os problemas propostos.
\item
  É capaz de produzir soluções, em formato de texto, compreensíveis, com
  nível de detalhamento e rigor adequados para que outros profissionais
  familiarizados com os conceitos possam entendê-las.
\item
  Aprendeu os conceitos e técnicas ensinados em sala de aula.
\item
  Sabe determinar quando a solução proposta resolve de fato o problema.
\item
  Consegue desenvolver uma argumentação lógica-dedutiva para chegar-se à
  resposta final.
\end{itemize}

\subsection*{Critério de correção das
VAs}\label{crituxe9rio-de-correuxe7uxe3o-das-vas}
\addcontentsline{toc}{subsection}{Critério de correção das VAs}

Com base nos objetivos apresentados acima, os critério de correção
serão:

\begin{enumerate}
\def\labelenumi{\arabic{enumi}.}
\tightlist
\item
  Interpretação correta das questões.
\item
  As soluções devem ser escrita com clareza, organização, rigor e
  detalhamento.
\item
  Soluções desenvolvidas fora do conteúdo ensinado, mesmo que corretas,
  não serão consideradas.
\item
  Questão com várias soluções será anulada.
\item
  Resposta correta com solução errada será anulada.
\end{enumerate}

Além disso, durante às VAs, serão adotadas as seguinte medidas:

\begin{itemize}
\tightlist
\item
  Proibido compartilhar material.
\item
  Só é permitida a saída após entrega definitiva da VA, ou seja, não é
  permitido ir ao banheiro ou beber água. A VA tem duração de no máximo
  2 horas.
\item
  É permitida a \textbf{consulta a material} (livros, notas de aula,
  caderno, apostilas e etc.) e uso de dispositivo computacional, como
  calculadora científica, ambos \textbf{sem acesso a internet}.
\item
  Fraude detectada, mesmo depois, zera a nota.
\end{itemize}

\subsection*{Segunda Chamada}\label{segunda-chamada}
\addcontentsline{toc}{subsection}{Segunda Chamada}

O Art. 98 do
\href{http://www.uff.br/sites/default/files/001-2015_regulamento_do_curso_de_graduacao_0.pdf\#page=55}{Regulamento
dos Cursos de graduação} garante ao aluno o direito a uma avaliação de
\textbf{Segunda Chamada (VR)}, \textbf{sem a necessidade de
justificativa!} Portanto, aquele aluno que \textbf{não puder ou não
quiser} fazer uma das VAs, poderá faltar que seu direito à Segunda
Chamada será garantido.

O conteúdo da \textbf{Segunda Chamada} será \textbf{toda a matéria do
semestre.}

\subsection*{Verificação
Suplementar}\label{verificauxe7uxe3o-suplementar}
\addcontentsline{toc}{subsection}{Verificação Suplementar}

Em concordância com o
\href{http://www.uff.br/sites/default/files/001-2015_regulamento_do_curso_de_graduacao_0.pdf\#page=55}{Regulamento
dos Cursos de graduação} em seu Art. 99, a \textbf{Verificação
Suplementar (VS)} é vetada aos discentes já aprovados e é
\textbf{obrigatória} para aqueles que tenham obtido média parcial entre
4.0 e 5.9, sendo esses dois limites incluídos.

\section{Listas de Exercícios}\label{listas-de-exercuxedcios}

Abaixo seguem os execícios da referência {[}1{]} que devem ser feitos. A
distribuição está dada por seção.

\begin{itemize}
\tightlist
\item
  \textbf{Parametrização, coordenadas polares, superfícies e funções
  vetoriais}\\
\end{itemize}

\textbf{§ 12.5 (Retas e Planos):} 3, 5, 7, 9, 23, 26, 31, 43, 45.\\
\textbf{§ 12.6 (Cilindros e Quádricas):} 3, 5, 13, 15, 21-28, 29, 31,
33.\\
\textbf{§ 10.3 (Coordenadas Polares):} 1, 3, 7, 9, 11, 15, 17, 19, 21,
23, 25, 29, 31, 33, 35, 37, 39, 41, 56.\\
\textbf{§ 13.1 (Funções Vetoriais):} 7, 9, 11, 19-24, 27, 37, 38, 39.\\
\textbf{§ 13.2 (Derivadas):} 3, 5, 17, 19, 21, 23, 25, 27 e 29.\\
\textbf{§13.3 (Comprimento de arco e Curvatura):} 1, 3, 5, 11, 13, 15,
17(b), 19(b), 21, 23.\\
\textbf{§ 13.4 (Velociade e Aceleração):} 3, 5, 11, 13, 19, 31,32, 41.\\
---

\begin{itemize}
\tightlist
\item
  \textbf{Funções de Várias Variáveis}\\
\end{itemize}

\textbf{§ 14.1 (Funções):} 11, 13, 15, 17, 19, 23, 25, 27, 30, 39, 41,
45, 49, 61, 63.\\
\textbf{§ 14.2 (Limite e Continuidade):} 5, 7, 9, 11, 13, 15, 29, 31,
35, 39, 41.\\
\textbf{§ 14.3 (Derivadas Parciais):} 15, 17, 21, 23, 25, 29, 31, 39,
41, 51, 53, 61, 63, 65.\\
\textbf{§ 14.4 (Aproximação Linear):} 1, 3, 5, 11, 13, 21, 25, 27, 33,
35, 37.\\
\textbf{§ 14.5 (Regra da Cadeia):} 1, 3, 5, 7, 15, 17, 19, 21, 23, 35,
43.\\
\textbf{§ 14.6 (Vetor Gradiente):} 5, 7, 9, 11, 15, 17, 21, 23, 25, 29,
31, 34, 39, 41, 53.\\
\textbf{§ 14.7 (Máximos e Mínimos):} 5, 7, 13, 15, 31, 33, 34, 35, 39,
41, 43, 47, 53.\\
\textbf{§ 14.8 (Multiplicadores de Lagrange):} 3, 5, 9, 15, 17, 19, 41,
43.

Exercícios do \textbf{Stewart Vol. 2, 7ª Edição.}

\textbf{§ 14.7:} 5, 7, 11, 13, 15, 31, 33, 34, 35, 39, 41, 43, 47, 53.\\
\textbf{§ 14.8:} 3, 5, 9, 15, 17, 19, ,21, 45.

\section{Referências}\label{sec-ref}

\protect\phantomsection\label{refs}
\begin{CSLReferences}{0}{0}
\bibitem[\citeproctext]{ref-stewart2}
\CSLLeftMargin{1. }%
\CSLRightInline{Stewart J (2011) C{á}lculo, volume 2, 6a edi{ç}{ã}o, São
Paulo, Editora Cengage Learning.}

\bibitem[\citeproctext]{ref-thomas2}
\CSLLeftMargin{2. }%
\CSLRightInline{Thomas GB (2011) C{á}lculo, volume 2, São Paulo, Editora
Addison Wesley.}

\bibitem[\citeproctext]{ref-diomara}
\CSLLeftMargin{3. }%
\CSLRightInline{Morgado MCF, Pinto D (2009) Cálculo diferencial e
integral de funções de várias variáveis, Rio de Janeiro, UFRJ.}

\bibitem[\citeproctext]{ref-tromba}
\CSLLeftMargin{4. }%
\CSLRightInline{Marsden JE, Tromba A (2003) Vector calculus, Rio de
Janeiro, Macmillan.}

\end{CSLReferences}

\section{Cronograma Proposto
Incialmente}\label{cronograma-proposto-incialmente}

Este foi o programa proposto no início do perído, que será mantido aqui
para fins de comparação com o cronograma que realmente foi efetivado
acima.

{ \textbf{1. Funções vetoriais e Superfícies} }

01 . ~ Seg -- 03/08 -- { Apresentação do estudantes. Apresentação do
curso. Aplicações. }\\
02 . ~ Ter -- 04/08 -- { Revisão de equação paramétrica da reta e
equação cartesiana do plano em \(\mathbb{R}^3\). Coordenadas Polares:
definição. }\\
{🏖️} { ~ Seg -- 10/08 -- Banca de Concurso }\\
{🏖️} { ~ Ter -- 11/08 -- Banca de Concurso }\\
03 . ~ Seg -- 17/08 -- { Curvas em Coordenadas Polares. Superfícies
Cilíndricas. Superfícies Quádricas. }\\
04 . ~ Ter -- 18/08 -- { Exemplos de Superfícies Quádricas. }\\
05 . ~ Seg -- 24/08 -- { Funções Vetoriais. Limite e continuidade.
Curvas parametrizadas no Plano. Curvas parametrizadas no Espaço.
Hélices. }\\
06 . ~ Ter -- 25/08 -- { Parametrização de curvas como interseção de
duas superfícies. Derivadas e regras de derivação de funções vetoriais.
}\\
07 . ~ Seg -- 31/08 -- { Comprimento de arco. Parametrização pelo
comprimento de arco. Vetor Tangente. }\\
08 . ~ Ter -- 01/09 -- { Definição de Curvatura. Vetor Normal e
Binormal. }\\
{🏖️} { ~ Seg -- 07/09 -- Independência }\\
09 . ~ Ter -- 08/09 -- { Movimento, velocidade e aceleração. }\\
10 . ~ Seg -- 14/09 -- { Fórmulas da curvatura }\\
11 . ~ Ter -- 15/09 -- { Aula de Exercícios 1 }\\
{ 12 . ~ Seg -- 21/09 -- VA 1 (Verificação de Aprendizagem) }\\
---

{ \textbf{2. Funções de Várias Variáveis} }

13 . ~ Ter -- 22/09 -- { Definição de Funções de Várias Variáveis.
Gráficos de funções de 2 variáveis. Curvas de Nível. Superfícies de
nível. }\\
14 . ~ Seg -- 28/09 -- { Limites e continuidade. Exemplos. }\\
15 . ~ Ter -- 29/09 -- { Derivadas Parciais. Derivadas de ordem
superior.Função de classe \(C^1\). }\\
16 . ~ Seg -- 05/10 -- { Regra da cadeia. Vetor Gradiente. Derivadas
Direcionais. }\\
17 . ~ Ter -- 06/10 -- { Derivadas direcionais. Propriedades do vetor
Gradiente }\\
{🏖️} { ~ Seg -- 12/10 -- Nossa Senhora Aparecida }\\
18 . ~ Ter -- 13/10 -- { Valores Máximo e Mínimo. Pontos críticos. Teste
da Derivada Segunda }\\
{🏖️} { ~ Qui -- 15/10 -- Dia do Professor }\\
19 . ~ Seg -- 19/10 -- { Topologia do \(\mathbb{R}^2\). Máximos e
Mínimos absolutos. Teorema de Weiestrass. }\\
20 . ~ Ter -- 20/10 -- { Multiplicadores de Lagrange em
\(\mathbb{R}^2\). Multiplicadores de Lagrange em \(\mathbb{R}^3\) }\\
{🏖️} { ~ Seg -- 26/10 -- Semana Nacional de Ciência e Tecnologia }\\
{🏖️} { ~ Ter -- 27/10 -- Semana Nacional de Ciência e Tecnologia }\\
{🏖️} { ~ Qua -- 28/10 -- Servidor Público }\\
{🏖️} { ~ Seg -- 02/11 -- Finados }\\
21 . ~ Ter -- 03/11 -- { Plano tangente. Aproximação linear. }\\
22 . ~ Seg -- 09/11 -- { Diferenciais em \(\mathbb{R}^2\) e
\(\mathbb{R}^3\). }\\
23 . ~ Ter -- 10/11 -- { Derivação Implícita e Teorema da Função
Implícita em \(\mathbb{R}^2\). }\\
{🏖️} { ~ Dom -- 15/11 -- Proclamação da República }\\
24 . ~ Seg -- 16/11 -- { Derivação Implícita e Teorema da Função
Implícita em \(\mathbb{R}^3\). }\\
25 . ~ Ter -- 17/11 -- { Aula de Exercícios 3 }\\
{🏖️} { ~ Sex -- 20/11 -- Consciência Negra }\\
{ 26 . ~ Seg -- 23/11 -- VA 2 (Verificação de Aprendizagem) }\\
---

{ \textbf{Provas Finais} }

{ 27 . ~ Ter -- 24/11 -- VR }\\
28 . ~ Seg -- 30/11 -- { --- }\\
{ 29 . ~ Ter -- 01/12 -- VS }\\
30 . ~ Seg -- 07/12 -- { -- }\\
31 . ~ Ter -- 08/12 -- { -- }\\
{🔚} ~ Sex -- 11/12 -- Fim do Período. {🙌} ---




\end{document}
